\label{sec:unions-and-intersections}

\begin{topicbox}{Closure Laws for Open and Closed Sets}
\begin{exposition}
The open closure laws of \S\ref{sec:open-sets} dualize to closed sets by complementation. Collecting them gives the bookkeeping rules used throughout the rest of the development.
\end{exposition}
\end{topicbox}
\begin{theorembox}{Theorem (Closure Laws for Closed Sets)}
\begin{theorem}[Closure Laws for Closed Sets]
\label{thm:closed-set-closure-operations}
An arbitrary intersection of closed subsets of $\mathbb{R}$ is closed, and a finite union of closed subsets of $\mathbb{R}$ is closed.
\hyperref[prf:closed-set-closure-operations]{Go to proof.}
\end{theorem}
\end{theorembox}
\begin{remark*}[Standard quantified statement]
\[
\Bigl(\forall i\in I\;F_i\text{ closed}\Bigr)\Longrightarrow \Bigl(\bigcap_{i\in I} F_{i}\Bigr) \text{ closed};\qquad
\Bigl(\bigwedge_{k=1}^{n} F_k\text{ closed}\Bigr)\Longrightarrow \Bigl(\bigcup_{k=1}^{n} F_{k}\Bigr) \text{ closed}.
\]
\end{remark*}
\begin{remark*}[Interpretation]
The roles of union and intersection swap relative to open sets: closed sets are stable under arbitrary intersection but only finite union. The asymmetry mirrors the open case under De~Morgan's laws.
\end{remark*}
\begin{dependencies}
\begin{itemize}
  \item \hyperref[def:closed-set]{Closed Set in $\mathbb{R}$}
  \item \hyperref[thm:open-set-closure-operations]{Unions and Finite Intersections of Open Sets}
\end{itemize}
\end{dependencies}
